% !TEX program = xelatex
\documentclass[12pt,a4paper]{article}

\usepackage{fontspec}
\usepackage{polyglossia}
\setmainlanguage{english}
\setmainfont{DejaVu Serif}
\setsansfont{DejaVu Sans}
\setmonofont{DejaVu Sans Mono}

\usepackage{amsmath,amssymb,amsfonts,mathtools}
\usepackage{bm}
\usepackage{geometry}
\usepackage{enumitem}
\usepackage{booktabs}
\usepackage{array}
\usepackage{hyperref}
\usepackage{microtype}
\geometry{margin=2.5cm}
\hypersetup{colorlinks=true,linkcolor=black,urlcolor=blue,citecolor=black}
\setlength{\emergencystretch}{3em}
\sloppy

\newcommand{\statementbox}[1]{%
\begin{center}%
\fbox{\begin{minipage}{0.88\linewidth}\centering #1\end{minipage}}%
\end{center}%
}

\DeclareMathOperator{\Adm}{Adm}
\DeclareMathOperator{\Struct}{Struct}
\DeclareMathOperator{\Dist}{Dist}
\DeclareMathOperator{\Tr}{Tr}
\DeclareMathOperator{\Diff}{Diff}
\DeclareMathOperator{\Var}{Var}
\DeclareMathOperator{\SE}{SE}
\DeclareMathOperator{\CI}{CI}

\newcommand{\TMI}{\mathcal T_{\mathrm{MI}}}
\newcommand{\TMIRQG}{\mathcal T_{\mathrm{MI\text{-}QG}}}
\newcommand{\M}{\mathcal M_U}
\newcommand{\Gspace}{\mathfrak G}
\newcommand{\Hilb}{\mathcal H}
\newcommand{\Prob}{\mathbb P}
\newcommand{\Lag}{\mathcal L}
\newcommand{\Eff}{\mathrm{eff}}
\newcommand{\std}{\mathrm{std}}
\newcommand{\obs}{\mathrm{obs}}
\newcommand{\diag}{\mathrm{diag}}
\newcommand{\MIphys}{\mathcal M_I^{\mathrm{phys}}}
\newcommand{\Unc}{\mathcal U}

\title{\textbf{Theory of Manifold Interfaces}\\
\large TMI-QG-2.0: A Minimal Predictive Model\\
\large of Interface-Induced Geometric Decoherence}
\author{Working publication-oriented TEX version}
\date{Version: 2.0 draft}

\begin{document}
\maketitle

\begin{abstract}
This text assembles the minimal predictive version of the quantum-gravitational layer of the Theory of Manifold Interfaces. The model does not claim to be a completed theory of quantum gravity and does not derive the Standard Model. Its task is more modest and more precise: to formulate an effective interface-field mechanism of geometric decoherence, in which the distinguishability of geometric alternatives generates an additional contribution to the decoherence rate. The central prediction is
\[
\Delta\Gamma_I
=
\Gamma_{\obs}-\Gamma_{\std}
=
\kappa_I |c_1|^2 |c_2|^2
\ln^2\!\left(\frac{a_1}{a_2}\right).
\]
Thus TMI-QG-2.0 defines a testable hypothesis: if an interface channel of geometric decoherence exists, the observed excess over standard decoherence should scale with the interface geometric contrast
\[
x_I=|c_1|^2|c_2|^2\ln^2\!\left(\frac{a_1}{a_2}\right).
\]
If such dependence is absent, the minimal version of the model is rejected or requires modification.
\end{abstract}

\tableofcontents

\section{Status of the version}

\statementbox{TMI-QG-2.0 is not a final Theory of Everything and not a completed theory of quantum gravity. It is a minimal effective model of interface-induced geometric decoherence.}

The text fixes the transition from the general architecture of the Theory of Manifold Interfaces to its first testable model. The central physical chain is
\[
g
\Rightarrow
\Adm_c(g)
\Rightarrow
\Hilb_I(g)
\Rightarrow
\Psi_I
\Rightarrow
\Dist_{\mathrm{poss}}(\Psi_I;g)
\Rightarrow
\Struct_I.
\]
Its meaning is
\[
\boxed{\text{geometry defines admissibility;}}
\]
\[
\boxed{\text{the quantum state defines the distribution of possibilities;}}
\]
\[
\boxed{\text{the interface transforms the distribution of possibilities into structure.}}
\]

In version 2.0 this chain is compressed into a predictive model:
\[
(a_1,a_2,c_1,c_2,\kappa_I)
\Rightarrow
\Gamma_I
\Rightarrow
C_{\mathrm{TMI}}(t)
\Rightarrow
\Struct_I(t)
\Rightarrow
\text{statistical test}.
\]

\section{Axiomatic core of the minimal model}

\subsection{Axiom of geometric admissibility}

Let \(g_{\mu\nu}\) be a metric on spacetime \(\M\). The metric defines the causal structure and the set of admissible causal transitions:
\[
g
\Rightarrow
\mathcal C^+_{I,g}(x)
\Rightarrow
J_g^+(x)
\Rightarrow
\Adm_c(g).
\]
Here
\[
\mathcal C^{+}_{I,g}(x)
:=
\{v\in T_x\M\setminus\{0\}\mid g_x(v,v)\leq 0,\ v\ \text{is future-directed}\}
\]
is the interface representation of the future causal cone. Its boundary coincides with the light cone of standard Lorentzian geometry:
\[
\partial\mathcal C^{+}_{I,g}(x)
\equiv
\mathcal C^{+}_{\mathrm{light},g}(x).
\]

\subsection{Axiom of quantum distribution of possibilities}

For a fixed geometry \(g\), the interface space of possibilities is defined as
\[
\Hilb_I(g)=L^2\bigl(\Adm_c(g),\mathcal B_g,\mu_g\bigr),
\]
where \(\mathcal B_g\) is a chosen \(\sigma\)-algebra and \(\mu_g\) is an effective or regularized measure on the space of admissible transitions. A quantum state
\[
\Psi_I\in \Hilb_I(g)
\]
defines a distribution of possibilities:
\[
d\Prob_{\Psi,g}(\Phi)=|\psi_I(\Phi)|^2d\mu_g(\Phi).
\]

\subsection{Axiom of interface structuring}

The interface transforms a distribution of possibilities into structure:
\[
I_Q:
\Dist_{\mathrm{poss}}(\Psi_I;g)
\longrightarrow
\Struct_I.
\]
Operationally this may be written as a quantum channel:
\[
\mathcal E_I(\rho)=\sum_k K_k^{(I)}\rho K_k^{(I)\dagger},
\qquad
\sum_k K_k^{(I)\dagger}K_k^{(I)}=\mathbf 1.
\]
In the limiting case of complete reduction,
\[
\rho'_I\approx \sum_k p_k |s_k\rangle\langle s_k|,
\qquad
p_k=\Tr(E_k^{(I)}\rho),
\qquad
E_k^{(I)}=K_k^{(I)\dagger}K_k^{(I)}.
\]

\section{The minimal two-geometry model}

\subsection{Two geometric alternatives}

Let there be two admissible geometric configurations:
\[
|g_1\rangle,
\qquad
|g_2\rangle.
\]
They form the minimal space of alternatives:
\[
\Hilb_g^{(2)}=\operatorname{span}\{|g_1\rangle,|g_2\rangle\}.
\]
For the first model we assume orthonormality:
\[
\langle g_i|g_j\rangle=\delta_{ij}.
\]

The quantum state of geometry is
\[
|\Psi_g\rangle=c_1|g_1\rangle+c_2|g_2\rangle,
\qquad
|c_1|^2+|c_2|^2=1.
\]

The density matrix is
\[
\rho_g(0)=|\Psi_g\rangle\langle\Psi_g|
=
\begin{pmatrix}
|c_1|^2 & c_1\overline{c_2}\\
c_2\overline{c_1} & |c_2|^2
\end{pmatrix}.
\]
The diagonal elements define probabilities, while the off-diagonal elements \(\rho_{12}\) and \(\rho_{21}\) define geometric coherence.

\subsection{Interface-induced decoherence}

Introduce the coherence-preservation coefficient \(\eta_I(t)\):
\[
\rho_{12}(t)=\eta_I(t)\rho_{12}(0).
\]
If \(0\leq\eta_I(t)\leq1\), then
\[
\rho_g(t)=
\begin{pmatrix}
|c_1|^2 & \eta_I(t)c_1\overline{c_2}\\
\eta_I(t)c_2\overline{c_1} & |c_2|^2
\end{pmatrix}.
\]

The interface-induced decoherence rate is \(\Gamma_I(t)\geq0\):
\[
\eta_I(t)=\exp\left(-\int_0^t\Gamma_I(s)ds\right).
\]
For a constant rate,
\[
\eta_I(t)=e^{-\Gamma_I t},
\qquad
|\rho_{12}(t)|=e^{-\Gamma_I t}|\rho_{12}(0)|.
\]

\subsection{Degree of structuring}

The minimal degree of structuring is
\[
\Struct_I(t):=1-\frac{|\rho_{12}(t)|}{|\rho_{12}(0)|}.
\]
Therefore,
\[
\Struct_I(t)=1-\eta_I(t),
\]
and, for a constant rate,
\[
\boxed{
\Struct_I(t)=1-e^{-\Gamma_I t}.
}
\]

\section{Action of the interface field}

\subsection{Minimal system}

The minimal physical system is
\[
(g_{\mu\nu},\Psi_g,\chi_I),
\]
where \(\chi_I\) is an effective interface field. The total action is
\[
S_{\mathrm{TMI}}
=
S_{\mathrm{GR}}[g]
+
S_Q[\Psi_g,g]
+
S_I[\chi_I,g]
+
S_{\mathrm{int}}[\chi_I,\Psi_g,g].
\]

The gravitational part is
\[
S_{\mathrm{GR}}[g]
=
\frac{c^4}{16\pi G}
\int_{\M}R\sqrt{-g}\,d^4x.
\]

\subsection{Interface-field Lagrangian}

In the minimal version \(\chi_I:\M\to\mathbb R\) is a real scalar field. Its Lagrangian is
\[
\Lag_\chi
=
-
\frac12g^{\mu\nu}\nabla_\mu\chi_I\nabla_\nu\chi_I
-
\frac12m_I^2\chi_I^2
-
\frac{\beta_I}{4}\chi_I^4
-
\frac12\xi_I R\chi_I^2
+
J_I\chi_I.
\]
Here \(m_I\) is the scale of the interface field, \(\beta_I\) is self-interaction, \(\xi_I\) is curvature coupling, and \(J_I\) is the source of the interface field.

The equation of motion is
\[
\boxed{
\Box_g\chi_I
-
(m_I^2+\xi_I R)\chi_I
-
\beta_I\chi_I^3
=
-J_I.
}
\]

\subsection{Source of the interface field}

The interface contrast for two geometries is defined in a soft normalized form:
\[
D_g(\Psi_g)=|c_1||c_2|d_I(g_1,g_2).
\]
The source is
\[
J_I=\alpha_I D_g(\Psi_g).
\]

In the stationary weak-field approximation,
\[
\nabla_\mu\chi_I\approx0,
\qquad
R\approx0,
\qquad
\beta_I\approx0.
\]
Then
\[
\chi_I=\frac{J_I}{m_I^2}=\frac{\alpha_I}{m_I^2}D_g(\Psi_g).
\]
If
\[
\Gamma_I=\lambda_I\chi_I^2,
\]
then
\[
\boxed{
\Gamma_I
=
\frac{\lambda_I\alpha_I^2}{m_I^4}
|c_1|^2|c_2|^2d_I^2(g_1,g_2).
}
\]

\section{Interface distance between geometries}

\subsection{General effective definition}

Strictly speaking, one should compare not coordinate representations of metrics but equivalence classes of geometries:
\[
[g]=\{\varphi^*g\mid \varphi\in\Diff(\M)\}.
\]
The effective interface distance is
\[
\boxed{
d_I([g_1],[g_2])
=
\frac{1}{\mathcal N_I}
\inf_{\varphi\in\Diff(\M)}
\left[
\int_\Omega
w_I(x)
\left\|g_1-\varphi^*g_2\right\|_{g_0}^2
 dV_{g_0}
\right]^{1/2}.
}
\]
Here \(w_I(x)\geq0\) is an interface weight, \(g_0\) is a regularizing background structure, and \(\mathcal N_I\) is a normalization factor.

\subsection{Minisuperspace toy model}

For the first computable case, geometry is represented by the scale factor \(a\):
\[
g_1=g(a_1),
\qquad
 g_2=g(a_2).
\]
Define
\[
\boxed{
d_I(g(a_1),g(a_2))=\frac{|\ln(a_1/a_2)|}{\sigma_a}.
}
\]
Then
\[
\Gamma_I
=
\frac{\lambda_I\alpha_I^2}{m_I^4\sigma_a^2}
|c_1|^2|c_2|^2
\ln^2\left(\frac{a_1}{a_2}\right).
\]
Let
\[
\kappa_I=\frac{\lambda_I\alpha_I^2}{m_I^4\sigma_a^2}.
\]
The central rate of the model becomes
\[
\boxed{
\Gamma_I
=
\kappa_I |c_1|^2|c_2|^2
\ln^2\left(\frac{a_1}{a_2}\right).
}
\]

\section{Main prediction of the model}

The normalized coherence is
\[
C_I(t)=\frac{|\rho_{12}(t)|}{|\rho_{12}(0)|}.
\]
Thus
\[
\boxed{
C_I(t)=
\exp\left[-
\kappa_I |c_1|^2|c_2|^2
\ln^2\left(\frac{a_1}{a_2}\right)t
\right].
}
\]

For the maximal superposition \(|c_1|^2=|c_2|^2=1/2\),
\[
\boxed{
\Gamma_I=\frac{\kappa_I}{4}\ln^2\left(\frac{a_1}{a_2}\right),
\qquad
C_I(t)=\exp\left[-\frac{\kappa_I}{4}\ln^2\left(\frac{a_1}{a_2}\right)t\right].
}
\]

The characteristic interface-decoherence time is
\[
\boxed{
\tau_I=\frac{1}{\Gamma_I}
=
\frac{1}{\kappa_I |c_1|^2|c_2|^2\ln^2(a_1/a_2)}.
}
\]

\section{Numerical example}

Take a maximal superposition and \(\kappa_I=1\). Then
\[
\Gamma_I=\frac14\ln^2(a_1/a_2).
\]

\begin{center}
\begin{tabular}{cccc}
\toprule
\(a_1/a_2\) & \(\ln(a_1/a_2)\) & \(\Gamma_I\) & \(\tau_I=1/\Gamma_I\)\\
\midrule
2 & 0.693 & 0.120 & 8.32\\
10 & 2.303 & 1.326 & 0.754\\
100 & 4.605 & 5.302 & 0.189\\
\bottomrule
\end{tabular}
\end{center}

For \(C_I(t)=e^{-\Gamma_I t}\):
\begin{center}
\begin{tabular}{cccc}
\toprule
\(t\) & \(a_1/a_2=2\) & \(a_1/a_2=10\) & \(a_1/a_2=100\)\\
\midrule
0 & 1.000 & 1.000 & 1.000\\
0.5 & 0.942 & 0.515 & 0.071\\
1 & 0.887 & 0.265 & 0.005\\
2 & 0.786 & 0.071 & 0.000025\\
5 & 0.548 & 0.0013 & \(\approx0\)\\
\bottomrule
\end{tabular}
\end{center}

The conclusion is
\[
\boxed{
\text{interface-induced decoherence grows as the square of the logarithm of the ratio of geometric scales.}
}
\]

\section{Comparison with standard decoherence}

Let standard decoherence have rate \(\Gamma_{\std}\):
\[
C_{\std}(t)=e^{-\Gamma_{\std}t}.
\]
TMI adds an interface rate:
\[
\Gamma_{\mathrm{total}}=\Gamma_{\std}+\Gamma_I.
\]
Therefore,
\[
C_{\mathrm{TMI}}(t)=e^{-(\Gamma_{\std}+\Gamma_I)t}.
\]
The relative difference is
\[
\frac{C_{\mathrm{TMI}}(t)}{C_{\std}(t)}=e^{-\Gamma_I t}.
\]
The excess suppression of coherence is
\[
\boxed{
\Delta_I(t)=1-e^{-\Gamma_I t}.
}
\]

In the minisuperspace model,
\[
\boxed{
\Delta_I(t)=1-
\exp\left[-
\kappa_I |c_1|^2|c_2|^2
\ln^2\left(\frac{a_1}{a_2}\right)t
\right].
}
\]

\section{Observable quantity and hypotheses}

The main observable excess is
\[
\Delta\Gamma_I=\Gamma_{\obs}-\Gamma_{\std}.
\]
The interface parameter of geometric contrast is
\[
\boxed{
x_I:=|c_1|^2|c_2|^2
\ln^2\left(\frac{a_1}{a_2}\right).
}
\]
Therefore, the TMI-QG prediction is
\[
\boxed{
\Delta\Gamma_I=\kappa_I x_I.
}
\]

The null hypothesis is
\[
\boxed{H_0:\ \kappa_I=0.}
\]
The TMI hypothesis is
\[
\boxed{H_{\mathrm{TMI}}:\ \kappa_I>0.}
\]
Equivalently,
\[
\boxed{
H_{\mathrm{TMI}}:
\Gamma_{\obs}
=
\Gamma_{\std}
+
\kappa_I |c_1|^2|c_2|^2
\ln^2\left(\frac{a_1}{a_2}\right).
}
\]

The minimal version is falsified if, with sufficient precision,
\[
\Gamma_{\obs}-\Gamma_{\std}=0
\]
for significant \(a_1/a_2\neq1\), or if the dependence does not scale as \(x_I\).

\section{Synthetic verification dataset}

Take
\[
|c_1|^2=|c_2|^2=\frac12,
\qquad
\kappa_I=0.5,
\qquad
\Gamma_{\std}=0.1.
\]
Then
\[
x_I=\frac14\ln^2(a_1/a_2),
\qquad
\Delta\Gamma_I=0.5x_I,
\qquad
\Gamma_{\obs}=0.1+0.5x_I.
\]

\begin{center}
\begin{tabular}{ccccc}
\toprule
\(a_1/a_2\) & \(\ln(a_1/a_2)\) & \(x_I\) & \(\Delta\Gamma_I\) & \(\Gamma_{\obs}\)\\
\midrule
1 & 0.000 & 0.000 & 0.000 & 0.100\\
2 & 0.693 & 0.120 & 0.060 & 0.160\\
5 & 1.609 & 0.648 & 0.324 & 0.424\\
10 & 2.303 & 1.325 & 0.663 & 0.763\\
50 & 3.912 & 3.826 & 1.913 & 2.013\\
100 & 4.605 & 5.302 & 2.651 & 2.751\\
\bottomrule
\end{tabular}
\end{center}

Coherence at \(t=1\):
\begin{center}
\begin{tabular}{cccc}
\toprule
\(a_1/a_2\) & \(C_{\std}(1)\) & \(C_{\mathrm{TMI}}(1)\) & \(\Delta_I(1)\)\\
\midrule
1 & 0.905 & 0.905 & 0.000\\
2 & 0.905 & 0.852 & 0.058\\
5 & 0.905 & 0.655 & 0.277\\
10 & 0.905 & 0.466 & 0.485\\
50 & 0.905 & 0.134 & 0.852\\
100 & 0.905 & 0.064 & 0.929\\
\bottomrule
\end{tabular}
\end{center}

\section{Statistical test form}

For scenarios \(i=1,\dots,n\), define
\[
x_I^{(i)}=|c_1^{(i)}|^2|c_2^{(i)}|^2
\ln^2\left(\frac{a_1^{(i)}}{a_2^{(i)}}\right),
\]
\[
y_i=\Delta\Gamma^{(i)}=\Gamma_{\obs}^{(i)}-\Gamma_{\std}^{(i)}.
\]
The minimal statistical model is
\[
\boxed{
y_i=\kappa_I x_I^{(i)}+\varepsilon_i,
\qquad
\mathbb E[\varepsilon_i]=0.
}
\]

If the errors have variances \(\sigma_i^2\), use weights
\[
w_i=\frac{1}{\sigma_i^2}.
\]
The estimator is
\[
\boxed{
\hat\kappa_I=
\frac{\sum_{i=1}^n w_i x_I^{(i)}y_i}{\sum_{i=1}^n w_i(x_I^{(i)})^2}.
}
\]
The standard error is
\[
\boxed{
\SE(\hat\kappa_I)=
\left(\sum_{i=1}^n w_i(x_I^{(i)})^2\right)^{-1/2}.
}
\]
The criterion for support is
\[
\boxed{
\hat\kappa_I>0
\quad\text{and}\quad
0\notin \CI(\hat\kappa_I).
}
\]
For an approximate 95\% criterion,
\[
\boxed{
\hat\kappa_I-1.96\,\SE(\hat\kappa_I)>0.
}
\]

\section{Domain of validity}

\subsection{What the model already provides}

The model already provides:
\begin{enumerate}[label=\arabic*.]
\item a computable quantity \(\Gamma_I\);
\item the coherence-suppression law \(|\rho_{12}(t)|=e^{-\Gamma_I t}|\rho_{12}(0)|\);
\item a distinction from standard decoherence;
\item a falsifiable hypothesis \(\Delta\Gamma_I\propto x_I\);
\item a statistical test form.
\end{enumerate}

\subsection{What the model does not yet prove}

The model does not prove:
\begin{enumerate}[label=\arabic*.]
\item a complete theory of quantum gravity;
\item a derivation of the Standard Model \(SU(3)\times SU(2)\times U(1)\);
\item the fundamental origin of all masses, charges, and particle generations;
\item a full measure \(\mu_G\) on the space of geometries;
\item a general distance formula between arbitrary geometries;
\item a universal mechanism of wave-function collapse.
\end{enumerate}

\subsection{Phenomenological parameters}

The coefficient \(\kappa_I\) is an effective parameter in the current version:
\[
\boxed{
\kappa_I\ \text{has not yet been derived from fundamental constants.}
}
\]
A stricter version should obtain
\[
\kappa_I=f(\lambda_I,\alpha_I,m_I,\sigma_a,\ldots)
\]
or more deeply,
\[
\kappa_I=f(S_I,\chi_I,g,\Psi).
\]

\subsection{The measure problem}

The measure \(\mu_G\) on the space of geometries is assumed to be effective or regularized in the current version. Possible future choices include:
\begin{enumerate}[label=\arabic*.]
\item a discrete measure;
\item a minisuperspace measure;
\item a cylindrical measure;
\item an effective path-integral measure.
\end{enumerate}

\section{Final formulation of TMI-QG-2.0}

\statementbox{The minimal TMI-QG model predicts an additional geometry-dependent contribution to the decoherence rate, generated by interface distinguishability of geometric alternatives.}

Main formula:
\[
\boxed{
\Delta\Gamma_I
=
\Gamma_{\obs}-\Gamma_{\std}
=
\kappa_I |c_1|^2|c_2|^2
\ln^2\left(\frac{a_1}{a_2}\right).
}
\]
Main testable dependence:
\[
\boxed{
\Delta\Gamma_I=\kappa_I x_I,
\qquad
x_I=|c_1|^2|c_2|^2\ln^2(a_1/a_2).
}
\]
Main statistical hypothesis:
\[
\boxed{H_0:\kappa_I=0,
\qquad
H_{\mathrm{TMI}}:\kappa_I>0.}
\]

Correct scientific status:
\[
\boxed{
\text{TMI-QG proposes an interface-field mechanism for geometric decoherence.}
}
\]

Incorrect status:
\[
\boxed{
\text{TMI-QG should not be presented as a completed quantum gravity or Theory of Everything.}
}
\]

\section{Short publication abstract}

The minimal version of TMI-QG formulates an effective interface-field mechanism of geometric decoherence. In the two-geometry minisuperspace approximation, quantum geometry is represented by the state
\[
|\Psi_g\rangle=c_1|g(a_1)\rangle+c_2|g(a_2)\rangle.
\]
The interface geometric contrast is defined as
\[
x_I=|c_1|^2|c_2|^2\ln^2(a_1/a_2).
\]
The model predicts an additional contribution to the decoherence rate:
\[
\Delta\Gamma_I=\kappa_I x_I.
\]
Therefore,
\[
C_{\mathrm{TMI}}(t)
=
\exp[-(\Gamma_{\std}+\kappa_I x_I)t].
\]
The model is falsifiable: if the observed excess \(\Gamma_{\obs}-\Gamma_{\std}\) does not show a positive linear dependence on \(x_I\), the minimal version of the interface channel of geometric decoherence is rejected or requires modification.

\section{Next stage}

After TMI-QG-2.0, the natural next steps are:
\begin{enumerate}[label=\arabic*.]
\item construct a stricter measure \(\mu_G\), or explicitly fix the regularization;
\item move from two geometries to an \(N\)-geometry and continuous model;
\item derive \(\kappa_I\) from a more fundamental action \(S_I[\chi_I,g,\Psi]\);
\item build a numerical toy model with a concrete potential \(U_{\mathrm{mini}}(a)\);
\item compare predictions with Wheeler--DeWitt, standard decoherence, path-integral approaches, semiclassical gravity, and LQG;
\item separate the physical TMI-QG branch from the mathematical branch of the category of interfaces.
\end{enumerate}

\end{document}
